\documentclass{article}
\usepackage{geometry}


\title{\vspace{-3.5cm}CSC263 Writing Assignment}
\author{Mock Interview}
\date{\small Due March 5}


\begin{document}
\maketitle


	\vspace{-1cm}
\section*{A Take-Home Interview Exercise}
\small At the beginning of the course, we mentioned that many employers ask data structure and algorithm questions in job
interviews. This assignment replicates a method that employers sometimes use to assess candidates: the take-home
interview assignment. Note that these types of assessments are not without controversy, as they can be time-consuming for
the candidate. Our goal is to use this situation as an example of a realistic situation so that you can practise explaining your
data structure ideas to a professional audience.

To that end, suppose that you are applying to be a Software Developer or Research Assistant at Super-Awesome-Org. The
technical team lead at the organization sends you the following email:
\vspace{1cm}

\begin{tabular}{|l|}
\hline
\begin{minipage}{0.8\linewidth}
\setlength{\parskip}{\baselineskip}
	\vspace{.2cm}
Dear (Your name),

Thank you for your interest in Super-Awesome-Org!

Your resume looks impressive, and we would love to bring you in for an interview. But before we schedule an on-site, we
want you to attempt the attached task so that we can discuss your solutions during the interview.
The task is to implement a dynamic hash table that uses quadratic probing. By dynamic hash table, we mean a hash table that automatically doubles in size when it reaches 50\% capacity or above, and reduces in size
once it drops to 25\% capacity or below.

There is some Python 3 starter code attached to help you get started. The code contains an incomplete definition of
a HashTable class. In particular, this class has a self.array attribute, which is the address table. There are some
additional specifications in this file as well (e.g. that we never go below the initial capacity), and some example calls at
the bottom of the file. For this task, please do not use Python’s built in dictionary.

In addition, please include an analysis of the amortized runtime for each of the following two scenarios:
1. Start with a hash table with an initial capacity of 10, insert N elements, and then delete all N elements.
2. Start with a hash table with an initial capacity of 2N that already has N-1 elements inserted into it. Then do N
alternating deletions and insertions (i.e. insert an element, delete an element, insert an element, delete an element,
etc.).

We want to get a sense of how you think about this type of analysis problem, and how you communicate.
Let me know if you have any questions. Otherwise, please send me your completed code, and a written explanation that
describes how your code works by March 5th. Thank you!

Riley F.

Technical Team Lead

Super-Awesome-Org
	\vspace{.2cm}
\end{minipage} \\ \hline
\end{tabular}

\section*{\large What to submit} 
Submit a PDF file of your email response to this request in the file csc263\_interview.pdf, and your code in csc263\_interview.py, on Markus.

The email should contain a description of your solutions and a justification of your approach, the amortized analysis that
was requested, and a note that the file csc263\_interview.py is included as an attachment.

Your email tone should be friendly and professional. It does not need to be overly formal (for example, personal pronouns
are appropriate), but should follow the usual format of a professional email:
\begin{itemize}
	\item Address your reader at the top of the email (e.g. “Hello Riley” or “Dear Riley”), and sign your name at the bottom of
the email.
	\item Make it clear to your reader in the first 1-2 sentences why you are sending this email.
	\item Use paragraphs to logically separate your thoughts. Each paragraph should have a topic sentence that gives the reader an
idea of what the paragraph will be about. Topic sentences are important so that readers can find important information
quickly, and to be able to skim through your email.
	\item We recommend that you end with a call to action, and describe what you would like the reader to do after reading the
email. In this case, an appropriate call to action could be to let you know what the next steps are if the organization
is interested in moving forward.
	\item Finally, sign off on the email by including a signature block.
You will be evaluated on not just the correctness and the explanation of your algorithm, but also the quality of your writing.
\end{itemize}
\section*{Marking scheme}
The point breakdown will be as follows:
\begin{itemize}
	\item {\bf (6 points) Unit tests.} The HashTable methods insert, delete, and search will be tested. We will be looking for
the correct values of attributes array, size, and capacity. You must submit and earn points on the written
part of the assignment in order to earn unit test points.
	\item {\bf (6 points) Explanation and Justification of Correctness.} Clearly explain the approach you used in your imple-
mentation. In addition to providing a complete explanation as to how these operations work, please also justify your
decisions and show that your approach is correct. We are looking for whether you thought of possible corner cases and
handled them appropriately, and whether your explanation aligns with your code. Note that this explanation should
not simply repeat your code. Instead, summarize your approach and why it works.
\begin{itemize}
\item  Example of a bad explanation: “The code loops through each element of a list, and if there is a 42, it adds one to
the counter.”
\item Example of a good explanation: “The code counts the number of times 42 appears in the list”
\end{itemize}
	\item {\bf (6 points) Amortized runtime analysis.} You may use a combination of logical and mathematical reasoning in this
analysis: i.e. you can use math or words to describe your runtime analysis, as long as your reasoning is clear to the
reader.
\item {\bf (4 points) Writing: Structure.} The email in your pdf file should follow the structure we outlined for a professional
email. Ideas should be introduced in a logical order, and the structure should be made clear to the reader. Use
paragraphs and sentences, with good transition words between them to guide the reader and let them know how your
sentences fit together. Use summaries and topic sentences to help readers anticipate what each paragraph will be about.
If a reader is looking for something specific (e.g. how you handle one specific case), they should be able to skim your
document and find exactly where you discuss that topic, without reading the entire document. You should mostly
use complete sentences, though you could also use bullets, tables and figures to highlight important information. An
observation should be presented before the analysis that relies on that observation is presented.
\item {\bf (4 points) Writing: Clarity/Writing Mechanics.} Any technical terms you introduce should be defined. Any pronouns
that you use (“this”, “it”, “they”) should be unambiguous—the reader should have no difficulty understanding what
pronouns refer to. Write concisely where possible, but not to the point of seeming robotic. For this assignment, please
limit your word count to be within 500. There should be few or no grammatical issues distracting to the readers. Opt
for simple sentences that introduce one idea, rather than complex sentences that link several ideas together. Keywords
(e.g. variable names) should be annotated with their “type”. Sentences should be structured to read well “left-to-right”,
with no new information that changes the meaning of the sentence introduced at the very end.
\begin{itemize}
\item Example of a sentence that doesn’t annotate keywords: “We have that p is not in q unless i is a prime.”
\item Example of a sentence that does annotate keywords: “We have that the number p is not in the list q unless its
index i is a prime.”
\item Example of a sentence that don’t read well “left-to-right”: “The code assumes that (42 in list) is not true.”
(The reader needs to rebuild their mental model of what you are saying when they read the last word in the
sentence.)
\item Example of a sentence that does read well “left-to-right”: “The code assumes that 42 is not in the variable list”
\end{itemize}
\item {\bf (4 points) Writing: Audience Expectations.} Your writing should be appropriately professional, and targeted towards a
potential employer, mentor, or collaborator. Avoid slang (e.g. “This code is lit”). 
 Use gender-neutral and inclusive
language. Assume that the reader is fairly technical (e.g. has a CS degree), but may not know the exact convention
that you are using. Avoid generic verbs like “do” or “make”
\begin{itemize}
\item Example of poor assumptions: “My analysis will assume all properties in slide 35 of lecture 14.”
\item Example of good assumptions: “My analysis will assume the universal hashing property.”
\item Example of generic term: “The code goes through the list and makes each item into a node.”
\item Example of technical verbs: “The code iterates through the list items, and constructs a node with each item.”
\end{itemize}
Additional Writing Advice:
\begin{itemize}
\item As always, start early!
\item Your first draft will suck, and that’s okay. Write the first draft to clarify your ideas and to get a deeper understanding
of the problem. Then, revise your writing so that it becomes more organized and more clear. Good writing always
takes several drafts; the first time through is often when you are just figuring out what exactly you want to say.
\item When you begin revising, keep the reader in mind. The secret to good writing is the ability to put yourself in the
reader’s shoes. What does the reader already know? What don’t they already know? What might they be expecting
to see next? Good writing shows that you care about the reader’s time!
\item When you write the email, you might be tempted to describe your thought process as you came up with the solution.
However, the reader only cares about your final result.
\item Use the above guidelines and the rubric as a checklist. Can you replace ambiguous pronouns with the corresponding
proper nouns? Do your sentences read well left-to-right? Can you replace generic verbs with technical verbs? (and so
on)
\end{itemize}
\end{itemize}
\end{document}
